\chapter*{Resumen}
\addcontentsline{toc}{chapter}{Resumen}

El conjunto de datos AI4Mars aporta máscaras de segmentación por píxel, producidas por anotación
colaborativa, para decenas de miles de imágenes tomadas en la superficie de Marte. Esas máscaras
se han empleado casi exclusivamente para entrenar redes neuronales. Este trabajo pregunta qué
indicadores cuantitativos de cobertura y organización de la roca pueden derivarse de ellas de
forma reproducible, y cuáles son sus límites de validez frente a diferencias en la anotación y
al juicio humano.

Se implementó un método clásico de procesamiento de imágenes que calcula, para cada
máscara, la cobertura de roca visible como cociente entre píxeles de roca y píxeles etiquetados,
y estima el número de rocas individuales mediante componentes conectadas, transformada de
distancia y división de aguas guiada por máximos de prominencia, con filtros explícitos de área
y forma. Se aplicó a las \cNEscenas{} imágenes de la cámara de navegación del rover
\textit{Curiosity} con máscara de entrenamiento.

La cobertura pudo calcularse en \cNCobCalc{} escenas (\cPctCobCalc); \cNCobPos{}
(\cPctCobPos{} del total) presentaron cobertura mayor que cero, con una mediana del \cCobMedVal{}
sobre píxeles etiquetados, y no se encontró evidencia de que sus valores altos se expliquen por
la fracción de escena etiquetada. El conteo, aplicado a las \cFlagOk{} escenas aptas, arrojó
\cNRocas{} rocas. Ambos indicadores presentan una asociación lineal prácticamente nula
($r = \cDepCCPearson$), aunque medidas no lineales detectan una dependencia débil. En las
\cNGold{} máscaras de especialista del propio conjunto, que recaen sobre otras imágenes, la
cobertura mediana de las escenas con roca es del \cGoldCobMedPos. Un segmentador entrenado,
usado como instrumento común sobre una muestra de cada fuente, atribuye el \cDescImgPct{} de la
diferencia de cobertura media entre las muestras a que las escenas de especialista son mucho
menos rocosas, y el \cDescAnotPct{} restante —unos dos puntos porcentuales, con un intervalo
que incluye el cero cuando se tiene en cuenta la agrupación de las imágenes por sol— a la fuente
de la anotación. Una evaluación exploratoria con un observador, que contó las rocas dentro de la región
anotada como roca grande, indica un subconteo sistemático: el procedimiento queda por debajo del
observador en \cVAbajo{} de \cVN{} escenas y por encima en \cVArriba{} (kappa ponderado de
\cVKappaPond). Cuenta de forma reproducible las instancias geométricamente separables dentro de
esa clase, pero no el número de rocas visibles: un mismo polígono agrupa a menudo varios bloques,
y la clase cubre solo una parte de la roca de la escena.

El resultado es, por tanto, negativo para el conteo: contar rocas individuales con estas
máscaras no es sencillo. Es, en cambio, diagnóstico para el conjunto de datos, cuyas
inconsistencias —escenas de especialista de composición muy distinta, una diferencia de
criterio entre anotaciones pequeña e incierta, una clase roca grande en la que ni los
especialistas coinciden y una anotación parcial de los bloques— se documentan con evidencia cuantitativa. El trabajo
entrega además un conjunto de resultados con una fila por imagen y veintiocho columnas, y un
diseño de evaluación que separa el efecto de la anotación del de las imágenes.

\vspace{0.4cm}
\noindent
\textbf{Palabras clave:} Marte; AI4Mars; segmentación semántica; división de aguas;
morfología matemática; cobertura de roca; anotación colaborativa; calidad de datos; validez de indicadores.
